% =====================================================================
%  CHƯƠNG 5 — CÁC GIẢI PHÁP VÀ ĐÓNG GÓP NỔI BẬT
%  Nguồn: datn_agent_skills-test/project_setup/architecture/DECISIONS.md
%  Mỗi mục theo khung 3 phần: (i) dẫn dắt vấn đề, (ii) giải pháp, (iii) kết quả.
%  KHÔNG lặp nội dung Chương 4 — chỉ tham chiếu chéo.
%  Số đo thực nghiệm chưa có nguồn đánh dấu  % TODO[cần đo]
% =====================================================================
\documentclass[../DoAn.tex]{subfiles}
\begin{document}

Chương này trình bày những giải pháp và đóng góp kỹ thuật nổi bật mà tác giả tâm
đắc nhất trong quá trình thực hiện đồ án. Mỗi đóng góp được đặt trong một mục độc
lập gồm ba phần: dẫn dắt vấn đề, giải pháp đề xuất, và kết quả đạt được. Các chi
tiết thiết kế đã mô tả ở Chương~4 không được lặp lại; mục này tập trung vào
\emph{vấn đề khó}, \emph{lập luận lựa chọn} và \emph{cách khắc phục}.

\section{Cầu nối CAN--Ethernet thời gian thực không mất khung}
\label{section:5.1}

\subsection*{(i) Vấn đề}
Nút cổng kết nối phải ghép hai môi trường có bản chất trái ngược: mạng CAN
\emph{hướng gói tin}, sự kiện rời rạc, độ ưu tiên thời gian thực cao; còn Ethernet
trên nền TCP lại \emph{hướng dòng byte}. Sự lệch pha tốc độ giữa hai phía dễ gây
mất khung CAN. Nghiêm trọng hơn, khi ứng dụng Android chưa mở cổng lắng nghe, lời
gọi \texttt{connect()} ở chế độ chặn (blocking) sẽ khoá tác vụ gửi cho tới khi
thiết lập được kết nối hoặc hết thời gian chờ; trong khoảng đó, ngắt nhận CAN vẫn
liên tục đẩy khung vào hàng đợi khiến hàng đợi tràn và khung bị loại bỏ ngay tại
ISR.

\subsection*{(ii) Giải pháp}
Tác giả thiết kế cơ chế đệm hai lớp kết hợp với socket \emph{không chặn}:
\begin{itemize}
  \item Ngắt nhận CAN chỉ làm một việc tối thiểu là đẩy khung vào hàng đợi an
    toàn-ngắt \texttt{canToEthQueue} rồi trả quyền điều khiển ngay, tách hoàn toàn
    phần thu nhận khỏi phần xử lý/truyền (xem~\ref{section:4.2.2}).
  \item Socket W5500 được mở với cờ \texttt{SF\_IO\_NONBLOCK}: \texttt{connect()}
    trả về \texttt{SOCK\_BUSY} tức thì, máy trạng thái tự tiến
    \texttt{INIT}$\to$\texttt{SYNSENT}$\to$\texttt{ESTABLISHED} qua các vòng lặp
    kế tiếp mà không khoá tác vụ gửi. Khi gặp timeout, chip tự đóng socket và mở
    lại sạch sẽ, tạo khả năng tự phục hồi.
\end{itemize}
Nhờ đó tác vụ gửi luôn rút được hàng đợi, độc lập với trạng thái sẵn sàng của
phía thu.

\subsection*{(iii) Kết quả}
Hiện tượng tràn hàng đợi và mất khung khi ứng dụng chưa sẵn sàng được loại bỏ;
socket tự tiến tới trạng thái thiết lập và tự khôi phục sau sự cố mạng mà không
cần can thiệp thủ công. Đóng góp này biến một liên kết "dễ gãy" thành một cầu nối
bền vững, là điều kiện tiên quyết cho trải nghiệm thời gian thực của cụm đồng hồ.
% TODO[cần đo]: thông lượng tối đa (frame/s) và tỉ lệ mất gói ở kịch bản tải cao (TC_ETH_02).

\section{Giao thức tải tin tối giản tận dụng đảm bảo của TCP}
\label{section:5.2}

\subsection*{(i) Vấn đề}
Một thiết kế giao thức theo lối kinh điển sẽ bọc dữ liệu trong khung có byte đồng
bộ, trường độ dài, mã loại thông điệp và mã kiểm lỗi (CRC). Cách làm này tăng chi
phí xử lý ở cả hai đầu và độ phức tạp phân giải, trong khi phần lớn các bảo đảm đó
đã được lớp vận chuyển TCP cung cấp sẵn. Câu hỏi đặt ra là: với một liên kết
TCP tin cậy, đâu là dạng tải tin tối giản nhưng vẫn đúng đắn?

\subsection*{(ii) Giải pháp}
Tác giả lựa chọn truyền \emph{nguyên khối} cấu trúc \texttt{VehicleState\_t} cố
định 18~byte: Gateway \texttt{memcpy} cấu trúc vào bộ đệm rồi \texttt{send()},
ứng dụng \texttt{readFully(18)} để gom đúng một ảnh chụp trạng thái. Do TCP đã bảo
đảm thứ tự và toàn vẹn dòng byte, hệ thống không bổ sung byte đồng bộ hay CRC ở
lớp ứng dụng; kích thước cố định cho phép phân giải bằng length-framing tầm
thường. Đi kèm là cơ chế \emph{tích luỹ lai}: một cấu trúc trạng thái bền vững
luôn giữ control mới nhất và sensor mới nhất, mỗi khung CAN hợp lệ kích hoạt gửi
trọn ảnh chụp đầy đủ (chi tiết bố cục byte ở~\ref{section:4.2.4},
Bảng~\ref{tab:tcp-packet}).

\subsection*{(iii) Kết quả}
Ứng dụng luôn nhận một trạng thái đầy đủ, nhất quán và không bao giờ phải hợp nhất
cập nhật từng phần. Mã phân giải ở phía Android được rút gọn tối đa, đồng thời
loại bỏ rủi ro lệch khung do phân tích sai byte đồng bộ. Giải pháp thể hiện
nguyên tắc \emph{không lặp lại bảo đảm đã có ở tầng dưới} — một lựa chọn thiết kế
đề cao sự đơn giản và tính đúng đắn.

\section{Độ tin cậy lớp ứng dụng trên CAN và tự phục hồi bus-off}
\label{section:5.3}

\subsection*{(i) Vấn đề}
Hai bài toán độc lập cùng xuất hiện trên liên kết CAN. Thứ nhất, cần phát hiện
hỏng dữ liệu, mất khung và trùng khung ở mức ứng dụng, độc lập với bộ đếm lỗi
phần cứng của bxCAN. Thứ hai, trong cấu hình ban đầu (one-shot, không bật
\texttt{AutoBusOff}), khi thiếu phản hồi ACK trên liên kết hai nút, mailbox truyền
bị kẹt và nút rơi vào trạng thái bus-off vĩnh viễn — biểu hiện là bộ đếm truyền
"đóng băng" và không tự phục hồi.

\subsection*{(ii) Giải pháp}
\begin{itemize}
  \item \textbf{Kiểm lỗi ở lớp ứng dụng:} mỗi khung mang một mã CRC8 (đa thức
    0x07) và một bộ đếm số thứ tự 4-bit. Phía Gateway tính lại CRC8 và đối chiếu,
    đồng thời theo dõi bước nhảy/trùng số thứ tự, ghi nhận vào bộ chẩn đoán
    \texttt{canDecodeDiag} (\texttt{crc\_error}, \texttt{seq\_gap},
    \texttt{seq\_duplicate}).
  \item \textbf{Cấu hình độ tin cậy đúng ngữ cảnh:} bật
    \texttt{AutoRetransmission} và \texttt{AutoBusOff} ở cả hai nút. Vì đây là
    liên kết telemetry hai nút (không phải bus kích hoạt theo thời gian), việc cho
    phép truyền lại và tự phục hồi bus-off là phù hợp.
\end{itemize}

\subsection*{(iii) Kết quả}
Hệ thống vừa phát hiện được lỗi đường truyền ở mức ứng dụng (qua các bộ đếm chẩn
đoán xuất ra UART), vừa tự thoát khỏi trạng thái bus-off thay vì treo vĩnh viễn.
Sự kết hợp giữa phát hiện lỗi (CRC/sequence) và phục hồi tự động (AutoBusOff)
mang lại một liên kết CAN vừa quan sát được vừa bền vững.
% TODO[cần đo]: số liệu crc_error/seq_gap thu được qua thời gian chạy thực (TC_CAN_03).

\section{Khắc phục lỗi đồng thời và khởi tạo trong hệ nhúng}
\label{section:5.4}

\subsection*{(i) Vấn đề}
Quá trình tích hợp bộc lộ những lỗi điển hình nhưng khó truy vết của lập trình
đồng thời và thứ tự khởi tạo:
\begin{itemize}
  \item \textbf{Tranh chấp khởi tạo (Gateway):} nếu bật ngắt nhận CAN trong
    \texttt{main()} \emph{trước} khi bộ lập lịch chạy, một khung CAN đến trong lúc
    \texttt{W5500\_Init()} đang thực thi sẽ khiến ISR gọi
    \texttt{xQueueSendFromISR} vào một hàng đợi còn NULL, kích hoạt
    \texttt{configASSERT} và treo vi điều khiển \emph{trước cả khi} in được dòng
    khởi động — UART câm hoàn toàn, rất khó chẩn đoán.
  \item \textbf{Tranh chấp cổng mạng (Android):} khi Activity sở hữu máy chủ TCP,
    mỗi lần Activity bị tái tạo lại gọi mở cổng trên một thể hiện mới trong khi
    thể hiện cũ chưa nhả cổng, gây lỗi \texttt{EADDRINUSE}.
\end{itemize}

\subsection*{(ii) Giải pháp}
\begin{itemize}
  \item Dời \texttt{HAL\_CAN\_Start} và \texttt{HAL\_CAN\_ActivateNotification}
    vào \emph{đầu tác vụ} \texttt{Eth\_SendTask} — tức sau khi hàng đợi đã được
    tạo và bộ lập lịch đã chạy — đồng thời thêm guard kiểm tra hàng đợi khác NULL
    trong callback ngắt như lớp phòng vệ thứ hai.
  \item Chuyển quyền sở hữu máy chủ TCP từ Activity sang một singleton phạm vi
    tiến trình do dịch vụ nền (foreground service, \texttt{START\_STICKY}) quản
    lý, để cổng \texttt{:5000} chỉ được bind đúng một lần cho cả tiến trình, độc
    lập với vòng đời giao diện (xem~\ref{section:4.2.3}).
\end{itemize}

\subsection*{(iii) Kết quả}
Vi điều khiển khởi động ổn định, không còn treo do tranh chấp ISR/hàng đợi; máy
chủ TCP không còn lỗi tranh chấp cổng khi giao diện được tái tạo. Đóng góp này
cho thấy việc tôn trọng \emph{thứ tự khởi tạo} và phân tách \emph{quyền sở hữu tài
nguyên} theo đúng vòng đời là chìa khoá cho độ ổn định của hệ thống nhúng đa
nhiệm.

\section{Chẩn đoán liên kết Ethernet điểm-điểm}
\label{section:5.5}

\subsection*{(i) Vấn đề}
Khi nối thẳng W5500 vào cổng \texttt{eth0} của Raspberry Pi, socket liên tục kẹt
ở trạng thái \texttt{SYNSENT} và log báo "socket not ready, drop frame". Nguyên
nhân là địa chỉ đích được cấu hình trùng với IP của giao diện Wi-Fi (\texttt{wlan0})
chứ không phải giao diện thật mang gói tin (\texttt{eth0}); ngoài ra nếu đặt hai
giao diện cùng dải địa chỉ, máy chủ sẽ có hai interface cùng mạng con gây nhập
nhằng ARP/định tuyến.

\subsection*{(ii) Giải pháp}
Tách liên kết Ethernet điểm-điểm sang một \emph{mạng con riêng}
\texttt{192.168.10.0/24}, độc lập hoàn toàn với dải Wi-Fi dùng cho gỡ lỗi; địa
chỉ đích của W5500 được đặt đúng bằng IP của \texttt{eth0}. Đồng thời bổ sung các
log chẩn đoán đọc trực tiếp thanh ghi chip (phiên bản chip, trạng thái liên kết
PHY, tốc độ/song công thực) để phân biệt nhanh ba nhóm sự cố: hỏng SPI/chip, đứt
liên kết vật lý, và sai địa chỉ đích (xem~\ref{section:4.1.2}).

\subsection*{(iii) Kết quả}
Socket tiến tới \texttt{ESTABLISHED} ổn định; các sự cố mạng được khoanh vùng
nhanh nhờ log chẩn đoán có chủ đích thay vì thông báo chung chung. Đây là một bài
học thực tiễn về việc \emph{địa chỉ đích phải khớp giao diện thật mang gói tin} —
một cái bẫy phổ biến trong hệ thống nhiều giao diện mạng.

\section{Toàn vẹn tín hiệu tương tự và chuẩn hoá thang đo đúng tầng}
\label{section:5.6}

\subsection*{(i) Vấn đề}
Hai vấn đề liên quan đến chuỗi tín hiệu tương tự. Thứ nhất, khi đọc ba biến trở
qua ADC ở chế độ quét, với thời gian lấy mẫu ngắn, tụ lấy-giữ của ADC không kịp
nạp giữa các kênh khiến kênh sau "ăn" giá trị kênh trước (rò điện tích). Thứ hai,
giá trị ADC 12-bit (0--4095) được truyền nguyên qua CAN/TCP, nhưng nếu phía hiển
thị chia cho toàn thang 16-bit (65535) thì kim đồng hồ chỉ chạy được khoảng
$1/16$ dải.

\subsection*{(ii) Giải pháp}
\begin{itemize}
  \item Tăng thời gian lấy mẫu ADC lên mức tối đa
    (\texttt{ADC\_SAMPLETIME\_239CYCLES\_5}) cho cả ba kênh, đủ để tụ lấy-giữ nạp
    đầy giữa các kênh đối với nguồn biến trở trở kháng cao (xem~\ref{section:4.2.1}).
  \item Sửa hệ số chuẩn hoá ở \emph{đúng tầng hiển thị} (Android), chia cho toàn
    thang ADC 12-bit \texttt{ADC\_FULL\_SCALE = 4095}, không can thiệp vào dây/CAN.
\end{itemize}

\subsection*{(iii) Kết quả}
Ba kênh ADC cho giá trị độc lập, ổn định; kim đồng hồ quét trọn dải biểu diễn
(vận tốc tới 240~km/h). Quan trọng hơn, việc khắc phục được thực hiện \emph{đúng
tầng trách nhiệm} (sửa thang đo ở tầng hiển thị thay vì sửa firmware), minh hoạ
nguyên tắc phân lớp được giữ vững xuyên suốt hệ thống.

\end{document}
